%% Taming divergent Knuth--Bendix completion with iterated functions
%% Prepared for submission to the Journal of Symbolic Computation (Elsevier).
\documentclass[preprint,12pt]{elsarticle}

\usepackage[utf8]{inputenc}
\usepackage[T1]{fontenc}
\usepackage{lmodern}
\usepackage{amsmath,amssymb,amsthm}
\usepackage{booktabs,array}
\usepackage{xcolor}
\usepackage{listings}
\usepackage{url}
\usepackage[hidelinks]{hyperref}

\newtheorem{definition}{Definition}
\newtheorem{proposition}{Proposition}
\newtheorem{example}{Example}

\newcommand{\R}{\mathcal{R}}
\newcommand{\E}{\mathcal{E}}
\newcommand{\rw}{\to}
\newcommand{\pw}[2]{#1^{#2}}               % power of a unary symbol
\newcommand{\pp}{\mathbin{\texttt{\^{}\^{}}}}     % HOL function power operator
\newcommand{\fp}[2]{(#1 \pp #2)}                  % HOL function power
\newcommand{\isa}[1]{\texttt{#1}}
\newcommand{\Suc}{\mathit{Suc}}
\newcommand{\sk}{\mathit{sk}}

\definecolor{kw}{RGB}{0,0,160}
\lstdefinelanguage{Isar}{
  morekeywords={theory,imports,begin,end,typedecl,consts,datatype,fun,where,
    lemma,by,apply,declare,complete,sorry,ML,thm,and},
  sensitive=true,
  morecomment=[n]{(*}{*)},
  morestring=[b]",
  literate={⇒}{{$\Rightarrow$}}1 {≡}{{$\equiv$}}1
    {‹}{{$\langle$}}1 {›}{{$\rangle$}}1
}
\lstset{language=Isar,basicstyle=\ttfamily\small,keywordstyle=\color{kw}\bfseries,
  columns=fullflexible,keepspaces=true,showstringspaces=false,upquote=true,
  frame=single,framesep=3pt,
  xleftmargin=4pt,xrightmargin=4pt,breaklines=true}

\journal{Journal of Symbolic Computation}

\begin{document}

\begin{frontmatter}

\title{Taming divergent Knuth--Bendix completion with iterated functions
in a higher-order theorem prover}

\author[aff]{Omar Monta\~{n}o-Rivas}
% TODO: complete the affiliation and the corresponding author's e-mail.
\address[aff]{Affiliation to be completed}

\begin{abstract}
Knuth--Bendix completion may run forever, generating an infinite family of
rewrite rules that differ only in the number of times a function is
iterated, as in $f(\pw{g}{n}(f(x))) \rw f(\pw{g}{n}(x))$. The classical
remedies schematize such families at the meta level, with meta-rules,
recurrence terms, terms with integer exponents or primal grammars, and then
need dedicated unification and rewriting algorithms for the schematic
objects. We observe that no new formalism is needed in higher-order logic,
where the whole family is a \emph{single theorem} about the function power
$g^n$ that can be proved by induction on $n$ inside the logic. We
present a \emph{divergence critic} for the theorem-producing completion
procedure of IsaScheme, a theory exploration system for Isabelle/HOL. The
critic recognizes families of chosen rules whose towers of unary symbols
grow linearly, generalizes them to a rule with a natural-number exponent,
proves it by induction with a proof method that the user can configure,
checks the termination of its first instances with
AProVE and restarts completion with the new rule, which is applied by a
simplification procedure that unfolds the right instance on demand. Every
rule of the final system is an Isabelle theorem, so soundness does not
depend on the critic. We describe the method and evaluate it on examples of
divergent completion taken from the literature, showing both the cases it
solves and the ones that are out of its reach.
\end{abstract}

\begin{keyword}
Knuth--Bendix completion \sep divergence \sep schematization \sep
function power \sep theory exploration \sep Isabelle/HOL
\end{keyword}

\end{frontmatter}

%% ===========================================================================
\section{Introduction}

Knuth--Bendix completion~\cite{KnuthBendix70} turns a set of equations
$\E$ into an equivalent convergent term rewriting system $\R$, which then
decides the word problem of $\E$ by comparing normal forms. A run of the
procedure may \emph{succeed} with a finite convergent system, \emph{fail}
because an equation cannot be oriented, or \emph{diverge} and keep adding
rules forever~\cite{Dershowitz89,BaaderNipkow98}. Unfailing
completion~\cite{BachmairDershowitzPlaisted89} removes the second outcome
by keeping unorientable equations and using them for ordered rewriting,
but divergence remains. It cannot be avoided in general, because some
finite theories with decidable word problem have no finite convergent
presentation at all~\cite{KapurNarendran85}, and the infinite
systems produced by divergent runs can be arbitrarily
complex~\cite{HermannPrivara86,SattlerKlein91}.

A small example shows the typical shape of divergence. The single rule
\[
  f(g(f(x))) \rw f(g(x))
\]
overlaps with itself, and the resulting critical pair produces
$f(g(g(f(x)))) \rw f(g(g(x)))$, which overlaps again, and so on. In the
limit, completion produces the infinite convergent system
\[
  \{\, f(\pw{g}{n}(f(x))) \rw f(\pw{g}{n}(x)) \mid n \ge 1 \,\}.
\]
The infinite system is regular, since its rules differ only in the length
of a tower of the unary symbol $g$. Since the work of H.~Kirchner on
meta-rules~\cite{Kirchner87,Kirchner89}, a whole line of research has
represented such families finitely, with recurrence
terms~\cite{ChenHsiangKong90}, terms with integer
exponents~\cite{Comon95}, primal grammars~\cite{HermannGalbavy97} or
enriched signatures~\cite{ThomasWatson93}, and has studied how to detect
them~\cite{Hermann89,KirchnerHermann90,ThomasJantke89,DershowitzPinchover90}.
In all these approaches the schematic rule lives at the \emph{meta level}.
Since it is not a rewrite rule of the object signature, it needs its own
matching, unification and critical pair computation, and its validity must
be argued about the infinite family outside the object logic.

This paper starts from a simple observation. In higher-order logic the
exponent is a first-class citizen. Isabelle/HOL~\cite{NipkowPaulsonWenzel02}
defines the function power $g \pp n$ by the two equations
$g \pp 0 = \mathit{id}$ and
$g \pp \Suc\;n = g \circ (g \pp n)$,
and the whole family above is the single HOL theorem
\[
  \forall n\,x.\ f(\fp{g}{\Suc\;n}\,(f\,x)) = f(\fp{g}{\Suc\;n}\,x),
\]
which can be proved by induction on $n$ from the initial rule. The
generalized rule is an ordinary theorem of the logic, so nothing needs to
be schematized at the meta level. Each of its instances is obtained by
instantiating $n$ and unfolding the definition of the function power.

We exploit this observation in the completion procedure of
IsaScheme~\cite{MontanoRivas12}, a theory exploration system for
Isabelle/HOL whose completion produces Isabelle theorems rather than
unchecked rewrite rules. Our contributions are:
\begin{itemize}
  \item a \emph{divergence critic} that monitors the rules chosen by
        completion, recognizes families whose towers of unary symbols grow
        linearly, and generalizes them into a single rule with a natural
        number exponent (Section~\ref{sec:critic});
  \item an automatic proof of the generalized rule by induction on the
        exponent, with a proof method that the user can replace by any Isar
        method, which keeps completion sound because every rule of the
        final system is a theorem checked by Isabelle's kernel;
  \item a way of \emph{using} such a rule during completion that leaves
        the termination checks, the term orderings and the computation of
        critical pairs unchanged. These components only ever see ordinary
        rules, while a simplification procedure unfolds, for each redex,
        the instance of the schematic rule that applies to it;
  \item an evaluation on divergent examples from the literature, including
        the cases that are out of reach of the method
        (Section~\ref{sec:examples}).
\end{itemize}

The paper is organized as follows. Section~\ref{sec:isascheme} gives a
short introduction to IsaScheme and to its completion procedure.
Section~\ref{sec:prelim} recalls the necessary notions of rewriting and
divergence. Section~\ref{sec:critic} presents the critic,
Section~\ref{sec:examples} the experiments, Section~\ref{sec:related}
related work and Section~\ref{sec:conclusion} our conclusions and future
work.

%% ===========================================================================
\section{IsaScheme in a nutshell}
\label{sec:isascheme}

IsaScheme~\cite{MontanoRivas12} is a system for \emph{theory exploration}
in Isabelle/HOL. Given a theory, for instance a few recursive functions
over an inductive datatype, it invents conjectures about them, discards the
false ones, proves the true ones and organizes the resulting theorems. Its
distinctive feature is the use of \emph{schemes}, higher-order formulas
that act as templates of interesting properties. The associativity scheme
\[
  \forall x\,y\,z.\ P\,(P\,x\,y)\,z = P\,x\,(P\,y\,z),
\]
for instance, yields one conjecture for every binary function $P$ of the
theory, and a scheme with an existentially quantified function variable
can be used to \emph{invent} a new function, which is the
concept-invention side of the system. Each conjecture is first tested with
counterexample generators (Quickcheck and Nitpick) and then given to an
automatic inductive prover. The prover analyses its own failures in the
style of proof critics~\cite{IrelandBundy96,Walsh96} and speculates the
missing lemmas.

The theorems discovered in this way are highly redundant, since an
associativity law produces infinitely many consequences, all of them
``interesting'' for a naive filter. IsaScheme therefore keeps the
discovered theorems as a \emph{convergent rewrite system} and runs
Knuth--Bendix completion every time a new theorem is added. A conjecture
is worth proving only if it is not already joinable, and the completed
system acts as a canonical, non-redundant summary of the theory explored so
far. Completion also lets the system replace definitions. Completing the
lemma $\mathit{itrev}\;xs\;[\,] = \mathit{rev}\;xs$, for instance, orients
it as $\mathit{rev}\;xs \rw \mathit{itrev}\;xs\;[\,]$ and deletes the
quadratic definition of $\mathit{rev}$ in favour of the linear one.

The current implementation of IsaScheme is the library IsaLibs, which we
have ported to Isabelle2026. It provides, among others, the Isar command
\isa{complete}, which is the entry point of everything that follows:
\begin{lstlisting}
complete "itrev xs [] = rev xs"
  by inductive_prover
\end{lstlisting}
The command opens a proof of the given equation. Once it is proved, the
equation is added to the defining equations of the functions it mentions
(and of the functions they depend on) and to the rules discovered before,
which are kept in the theorem collection \isa{compl}. The union is then
completed, and the trace reports the rules that were added to and removed
from \isa{compl}. The completion procedure has the following ingredients:
\begin{itemize}
  \item It is \emph{unfailing completion}~\cite{BachmairDershowitzPlaisted89},
        which keeps unorientable (permutative) equations, uses them for
        ordered rewriting and also computes critical pairs with them.
  \item Joinability of permutative equations is checked for all ground
        instances with the criterion of Martin and
        Nipkow~\cite{MartinNipkow90}, which tries all the ways of adding the
        variables of the equation to the precedence of a lexicographic path
        ordering.
  \item Instead of a fixed reduction ordering, orientation is decided by
        a \emph{termination oracle}, which accepts a candidate rule in a
        given direction when the termination prover
        AProVE~\cite{Giesl17,ArtsGiesl00} proves the termination of the
        current rules together with it. Two completions, which differ in
        the direction of the term ordering they use, run in parallel and
        the first one to finish is kept.
  \item Unification uses higher-order patterns~\cite{Miller91}, so the
        rules may contain $\lambda$-abstractions and function variables.
  \item Completion is \emph{theorem producing}, in the sense that every
        rule is an Isabelle theorem. The initial rule comes from the user's proof, critical
        pairs are derived by rewriting with theorems (Isabelle conversions),
        and all other inferences only decide which theorems are kept. As
        in any LCF-style system, errors in the termination oracle or in the
        heuristics may affect completeness or efficiency, but never
        soundness.
\end{itemize}
Since the system works on theorems, critical pair generation is optional
(option \isa{generate\_cps}) and off by default. In theory exploration,
completion is mainly used to interreduce and orient lemmas proved by
induction, and critical pairs are only needed when the user wants a
convergent system, which is precisely the situation in which completion
may diverge.

%% ===========================================================================
\section{Preliminaries}
\label{sec:prelim}

We use standard notation for term rewriting~\cite{BaaderNipkow98}. Terms
are built from a signature of typed function symbols and variables; we
write $s|_p$ for the subterm of $s$ at position $p$ and $s[t]_p$ for the
result of replacing it by $t$. A rewrite rule $l \rw r$ is an oriented
equation, and a term rewriting system (TRS) $\R$ is \emph{convergent} if
its rewrite relation is terminating and confluent. If two rules
$l_1 \rw r_1$ and $l_2 \rw r_2$ (with disjoint variables) are such that
$l_1|_p$ is not a variable and $\sigma$ is a most general unifier of
$l_1|_p$ and $l_2$, then
$\langle r_1\sigma,\ l_1\sigma[r_2\sigma]_p\rangle$ is a \emph{critical
pair}. A terminating TRS is confluent iff all its critical pairs are
joinable.

Completion~\cite{KnuthBendix70,BachmairDershowitzPlaisted89} can be
presented as an inference system on pairs $(\E,\R)$ of equations and
rules~\cite{BaaderNipkow98}, in which \emph{deduce} adds a critical pair to $\E$,
\emph{orient} moves an equation to $\R$ in a direction compatible with the
reduction ordering, \emph{simplify} and \emph{delete} normalize and remove
equations, and \emph{compose} and \emph{collapse} interreduce the rules. A
fair run that never fails produces in the limit a convergent system
$\R_\omega$ equivalent to the initial equations, which may be
infinite~\cite{HirokawaMiddeldorpSternagelWinkler17}. A run
\emph{diverges} if it never terminates; if the procedure is fair,
$\R_\omega$ is then infinite.

A function symbol $c$ is \emph{unary} if its type is
$\tau \Rightarrow \tau$ for a type $\tau$ that is not a function type. We
write $\pw{c}{k}(s)$ for the tower $c(c(\cdots c(s)\cdots))$ with $k$
occurrences of $c$. In Isabelle/HOL the function power is the constant
$\mathit{compow}$, written \isa{c \^{}\^{} n}, with the defining equations
\[
  c \pp 0 = \mathit{id}, \qquad
  c \pp \Suc\;n = c \circ (c \pp n),
\]
so that $\fp{c}{\Suc^k\;0}\,s$ unfolds to $\pw{c}{k}(s)$ by rewriting with
these equations and with the definitions of $\circ$ and $\mathit{id}$.

\begin{example}\label{ex:fg}
Let $\E = \{ f(g(f(x))) \approx f(g(x)) \}$, oriented from left to right.
The left-hand side overlaps with itself at position $1.1$ with the unifier
$x \mapsto g(f(y))$, so the term $f(g(f(g(f(y)))))$ rewrites at the root to
$f(g(g(f(y))))$ and at position $1.1$ to $f(g(f(g(y))))$, which rewrites
further to $f(g(g(y)))$. The critical pair
$f(g(g(f(y)))) \approx f(g(g(y)))$ is a new rule, which overlaps with the
first one in the same way, and completion generates
$f(\pw{g}{n}(f(x))) \rw f(\pw{g}{n}(x))$ for every $n \ge 1$.
\end{example}

%% ===========================================================================
\section{A divergence critic based on function powers}
\label{sec:critic}

The critic is integrated in the main loop of completion. After the
procedure \emph{chooses} a rule $\rho$ to compute its critical pairs, the
critic compares $\rho$ with the rules chosen before. If $\rho$ completes a
family, the critic tries to replace the whole family by a single rule with
an exponent. If it succeeds, completion is restarted with the new rule;
otherwise it continues as usual. The critic is only active when critical
pairs are generated.

\subsection{Tower signatures}

\begin{definition}[Tower skeleton and signature]
A \emph{maximal tower} in a term $t$ is a subterm $\pw{c}{k}(s)$, with $c$
unary and $k \ge 1$, that is not itself an argument of $c$ and such that $s$
is not headed by $c$. The \emph{skeleton} $\sk(t)$ is obtained by replacing,
from the root downwards, every maximal tower $\pw{c}{k}(s)$ by
$\Box_c(\sk(s))$, where $\Box_c$ is a hole labelled with the symbol $c$. The
\emph{tower lengths} $\ell(t) \in \mathbb{N}^m$ are the exponents $k$ of
those towers, listed in pre-order. The \emph{signature} of a rule
$l \rw r$ is the pair $(\sk(l \approx r), \ell(l \approx r))$, computed
after renaming the variables canonically in order of first occurrence.
\end{definition}

The rules of Example~\ref{ex:fg} all have the skeleton
$\Box_f(\Box_g(\Box_f(x))) \approx \Box_f(\Box_g(x))$, and the rule with
$n$ copies of $g$ has tower lengths $(1,n,1,1,n)$.

\subsection{Recognizing a family}

The critic keeps the history $H$ of the signatures of the chosen rules.
When a rule with signature $(S, v)$ is chosen, it looks for a previous
signature $(S, w) \in H$ with the same skeleton such that
\[
  d = v - w \in \{0,1\}^m, \qquad d \ne 0, \qquad (S, w - d) \in H.
\]
Three rules of the history are then consecutive members of the arithmetic
progression of tower lengths $u, u + d, u + 2d, \ldots$, where the base $u$
is obtained by going down from $w - d$ in steps of $d$ while the
corresponding signature is still in $H$. Requiring three members rather
than two avoids generalizing coincidences, and restricting the steps to
$\{0,1\}$ means that all the growing towers grow at the same rate. Each
skeleton is generalized at most once.

\subsection{Generalization and proof}

The candidate rule is rebuilt by filling the holes of the skeleton. A hole
$\Box_c$ whose tower does not grow ($d_j = 0$) receives the fixed tower
$\pw{c}{u_j}$, and a hole whose tower grows ($d_j = 1$) receives the
function power $c \pp (\Suc^{u_j}\;n)$, where $n$ is a
fresh natural number variable. The instance $n = i$ of the candidate is
then the member of the family with tower lengths $u + i\,d$. In
Example~\ref{ex:fg}, with $u = (1,1,1,1,1)$ and $d = (0,1,0,0,1)$, the
candidate is
\[
  f\,(\fp{g}{\Suc\;n}\,(f\,x)) = f\,(\fp{g}{\Suc\;n}\,x).
\]
The candidate is a conjecture, not a consequence of the rules seen so far,
so it must be proved. By default the critic applies induction on $n$, with
all the other free variables generalized, then simplifies both subgoals
with the current rules and finally calls the first-order prover
Metis~\cite{PaulsonSusanto07} with the same rules. The proof has a time
limit (option \isa{max\_time\_in\_critic\_proof}); if it fails, the
skeleton is marked as tried and completion goes on unchanged.

No fixed tactic proves every family, so the proof is configurable at two
levels. Lemmas declared with the attribute \isa{critic\_simps} are added
to the simplification and Metis steps of the default tactic. Alternatively,
the command
\begin{lstlisting}
divergence_critic_method ‹method›
\end{lstlisting}
replaces the default tactic by an arbitrary Isar method. The method text
is checked when it is declared and is stored in the theory context, so it
follows the usual scoping of Isabelle declarations. Since the names of the
variables of a conjecture are only known at run time, two auxiliary
notions make such methods writable. The proof method \isa{exponent\_induct}
performs induction on the unique natural-number variable that occurs in an
exponent of a function power, with all the other free variables
generalized. The fact \isa{critic\_rules} is bound, during the proof
attempt, to the current rules of the completion. In this notation, the
default tactic is close to
\begin{lstlisting}
exponent_induct; (simp add: critic_rules)?; metis critic_rules
\end{lstlisting}
Whatever the method, the result is checked by Isabelle's kernel, so the
choice of method affects which families are accepted, never soundness.

\subsection{Termination}

A rule with an exponent stands for infinitely many ordinary rules, one for
each value of the exponent, and the termination oracle can only check
finitely many of them. The critic
asks AProVE for the termination of the current rules together with the
instances $n = 0, 1, 2$ of the candidate. This check is a heuristic and
does not prove that the infinite system terminates. In the examples of
Section~\ref{sec:examples} it never rejected a proved rule; most of the
schematic rules found there remove an occurrence of a symbol, such as the
inner $f$ of Example~\ref{ex:fg}, at every instance. Proving termination of
the schematic system itself is left for future work
(Section~\ref{sec:conclusion}).

\subsection{Rewriting with a schematic rule}

A rule $\hat\rho(n)$ with an exponent cannot be given to the simplifier as
it is, because its left-hand side contains $g \pp n$,
which does not match a tower $\pw{g}{k}$, and because the rest of the
procedure (termination checks, term orderings, critical pairs) only works
on rules without exponents. Therefore schematic rules are kept in the proof context,
not in the rule lists, and are used through a \emph{simplification
procedure} (simproc). The simproc is triggered on terms headed by the head
symbol of the left-hand side; for a redex $t$, it tries the instances
$\hat\rho(0), \hat\rho(1), \ldots$ in turn, each obtained by instantiating
$n$ with a numeral $\Suc^i\;0$ and unfolding the function power, until one
of them rewrites $t$. The search is finite, because a tower in $t$ cannot be
longer than the size of $t$. Each instance is an instance of a theorem, so
the rewrite step is again justified by Isabelle's kernel, and the instances
are cached. The simproc is added to every simplifier set used by
completion, namely those of normalization, of the joinability tests and of
the final interreduction.

\subsection{Restarting completion}

When the critic succeeds, completion is restarted from the initial
equations, with the new rule in the context and with the critical pairs
among its instances $n = 0, 1, 2$ and between them and the initial
equations. Members of the family that are generated again are now
joinable, since the simproc rewrites them, and completion can terminate.
The final system consists of the rules computed by the restarted
completion and the schematic rules of the context, and it is added to
\isa{compl}. A schematic rule in \isa{compl}, or a schematic rule given
directly to \isa{complete}, is recognized by the presence of
$\mathit{compow}$ and put into the context in the same way. At most five
schematic rules are introduced in one completion, to avoid an unbounded
chain of generalizations.

\subsection{Soundness and completeness}

\begin{proposition}
Every rule of the system returned by \isa{complete}, including the
schematic ones, is a theorem of the current Isabelle theory.
\end{proposition}

\begin{proof}[Proof sketch]
The rules produced by completion without the critic are theorems, as
discussed in Section~\ref{sec:isascheme}. A schematic rule is accepted only
after a proof by induction checked by Isabelle's kernel, and its instances
and critical pairs are obtained from it by instantiation and rewriting with
theorems.
\end{proof}

Completeness is a different matter. The critical pairs of a schematic rule
are only approximated by those of its first three instances, and the
termination of the infinite family is only tested on those instances.
Hence the final system is not guaranteed to be convergent, even if every
check succeeds. A complete treatment would compute the critical pairs of
the schematic rule by unification of terms with integer
exponents~\cite{Comon95}; we discuss this in Section~\ref{sec:conclusion}.

%% ===========================================================================
\section{Examples from the literature}
\label{sec:examples}

We evaluated the critic on divergent completions reported in the
literature. Uninterpreted symbols are declared with \isa{typedecl} and
\isa{consts}, the equations are stated with \isa{axiomatization}, and each
equation is given to \isa{complete} with critical pairs enabled:
\begin{lstlisting}
typedecl T
consts f :: "T ⇒ T"  g :: "T ⇒ T"
axiomatization where fg: "f (g (f x)) = f (g x)"
declare [[generate_cps = true]]
complete "f (g (f x)) = f (g x)" by (rule fg)
\end{lstlisting}
Strings over an alphabet are encoded as compositions of unary functions,
so that the word $aba$ becomes the term $a(b(a(x)))$. All runs used
Isabelle2026 with 8 threads on a machine with two Intel Xeon E5645
processors and 30\,GB of memory, the AProVE server of February 2026, and
the default options of IsaLibs, which allow 20\,s for each termination
check and 60\,s for each proof attempted by the critic. Times are wall-clock times of the
\isa{complete} commands, measured from the trace. Runs without the critic
(option \isa{divergence\_critic = false}) were stopped after 300\,s; they
were run two at a time, so the number of rules they reached is only
indicative. The theories are distributed with the paper (see the
Availability section).
Table~\ref{tab:results} summarizes the results.

\begin{table}[t]
\centering
\small
\caption{Divergent completions from the literature. The column ``Without
critic'' runs completion with critical pairs and
\isa{divergence\_critic = false} and stops it after 300\,s. The column
``With critic'' uses the default configuration, except for row 3b. Times
are wall-clock times of the \isa{complete} commands.}
\label{tab:results}
\begin{tabular}{@{}l>{\raggedright}p{3.3cm}>{\raggedright}p{2.1cm}>{\raggedright}p{2.7cm}>{\raggedright\arraybackslash}p{3.3cm}@{}}
\toprule
& Equations & Source & Without critic & With critic\\
\midrule
1 & $f(g(f(x))) \approx f(g(x))$ & Example~\ref{ex:fg} &
  no end; members up to $n = 8$ &
  ends, 1 schematic rule, 3.8\,s\\
1b & $f(g(f(x))) \approx g(f(x))$ & variant of 1 &
  no end; members up to $n = 12$ &
  ends, 1 schematic rule, 0.4\,s\\
2 & Denzinger, IsaScheme's orientation & \cite[Fig.~48]{BurghardtHeinz96} &
  ends, 3 rules, 0.9\,s & same (critic not triggered)\\
2b & Denzinger, orientation of the LPO & \cite[Fig.~48]{BurghardtHeinz96} &
  no end & ends, 1 schematic rule, 139\,s; system not confluent\\
3 & $a(b(a(x))) \approx b(a(b(x)))$ &
  \cite{KapurNarendran85}, \cite[Ex.~1]{HirokawaMiddeldorpSternagelWinkler17} &
  not run (no finite convergent system exists) & family found but not proved; no end\\
3b & same, with \isa{funpow\_swap1} in \isa{critic\_simps} & & same as 3 &
  ends, 1 schematic rule, 2\,s; the convergent system of the literature\\
4 & associativity of Peano $+$ & \cite[\S4.4]{Dershowitz89} &
  no end; Dershowitz's sequence up to $n = 10$ & ends, 2 schematic rules, 73\,s\\
5 & bands & \cite[\S2.4]{Dershowitz89} & not run & not triggered; no end\\
6 & signed binary trees & \cite[Ex.~2.1]{Kirchner89} & not run & not triggered; no end\\
\bottomrule
\end{tabular}
\end{table}

\subsection{Self-overlapping rules}

For Example~\ref{ex:fg}, completion chooses the rules
$f(\pw{g}{k}(f(x))) \rw f(\pw{g}{k}(x))$ for $k = 1, 2, 3$ and, when the
third one is chosen, the critic reports:
\begin{lstlisting}[language={}]
Divergence critic: f ((g ^^ Suc n) (f x)) = f ((g ^^ Suc n) x)
Schematic rule: f ((g ^^ Suc ?n) (f ?x)) ≡ f ((g ^^ Suc ?n) ?x)
Restarting completion with schematic rule: ...
Done: f ((g ^^ Suc ?n) (f ?x)) ≡ f ((g ^^ Suc ?n) ?x)
\end{lstlisting}
The restarted completion finds that every critical pair is joinable with
the schematic rule, and the final system is the single schematic rule,
which subsumes the initial equation (its instance $n = 0$). The variant
$f(g(f(x))) \approx g(f(x))$ behaves in the same way and ends with
$f(\fp{g}{\Suc\;n}\,(f\,x)) \rw \fp{g}{\Suc\;n}\,(f\,x)$. Without the
critic, both completions keep generating the members of the family until
they are stopped; after 300\,s they had reached $n = 8$ and $n = 12$
respectively.

\subsection{Two families and the role of the ordering}

\begin{example}[Denzinger]\label{ex:denzinger}
Burghardt and Heinz~\cite[Fig.~48]{BurghardtHeinz96} report the following
example of J.~Denzinger. With the lexicographic path ordering for the
precedence $c \succ b \succ a \succ d \succ f \succ e$, completion of
\[
  a(c(x)) \approx x,\quad c(f(x)) \approx e(x),\quad
  a(e(x)) \approx f(x),\quad c(d(x)) \approx b(c(x))
\]
generates the two families $a(\pw{b}{n}(c(x))) \rw \pw{d}{n}(x)$ and
$a(\pw{b}{n}(e(x))) \rw \pw{d}{n}(f(x))$, for every
$n \ge 1$.\footnote{In the figure the last equation appears as the first
step of completion. Since $b$ and $d$ do not occur in the other equations,
it must be part of the input.}
\end{example}

Given these equations in this order, IsaScheme does not diverge at all.
Its termination oracle is not tied to a precedence, and it orients the
second and fourth equations from right to left, as $e(x) \rw c(f(x))$ and
$b(c(x)) \rw c(d(x))$. Completion then stops with the finite convergent
system
\[
  \{\, a(c(x)) \rw x,\ \ e(x) \rw c(f(x)),\ \ b(c(x)) \rw c(d(x)) \,\},
\]
whose left-hand sides do not overlap, and which proves the third equation
by $a(e(x)) \rw a(c(f(x))) \rw f(x)$. This is a reminder that divergence
is a property of a run, and depends on the
ordering~\cite{Lescanne86,Dershowitz89}.

To reproduce the divergent run, we fixed the orientations of the
precedence for the input equations with IsaScheme's command
\isa{orient\_rules}:
\begin{lstlisting}
orient_rules "c (f ?x) = e ?x"
orient_rules "c (d ?x) = b (c ?x)"
\end{lstlisting}
The rules derived during completion are still oriented by the oracle,
which prefers $d(x) \rw a(b(c(x)))$ and $f(x) \rw a(e(x))$, eliminating
$d$ and $f$. Completion now diverges, and the critic sees three families.
It first conjectures Denzinger's own family,
$a(\fp{b}{\Suc\;n}\,(c\,x)) = \fp{d}{\Suc\;n}\,x$. The restart, however,
comes from a second family, which it detects and proves 1.4\,s later:
\[
  c\,(a\,(\fp{b}{\Suc\;n}\,(c\,x))) = \fp{b}{\Suc\;n}\,(c\,x).
\] After the restart a third family appears,
$c(a(\pw{b}{n}(e(x)))) \rw \pw{b}{n}(e(x))$. The critic conjectures
$c\,(a\,(\fp{b}{\Suc\;n}\,(e\,x))) = \fp{b}{\Suc\;n}\,(e\,x)$ but does not
prove it within the time limit, and completion terminates after 139\,s.
The final system contains the schematic rule and the first three members
of the third family as separate rules.

This system is \emph{not} confluent. The terms
$c(a(\pw{b}{4}(e(x))))$ and $\pw{b}{4}(e(x))$ are equal, but both are
irreducible. The missing rule is the critical pair between the instance
$n = 3$ of the schematic rule and the rule $c(a(e(y))) \rw e(y)$, at the
position of the inner $c$, and the critic only computes the critical pairs
of the instances $n = 0, 1, 2$. The example shows concretely both
approximations discussed in Section~\ref{sec:critic}. It also shows what an
exact computation would buy. The critical pair of the schematic rule
itself with $c(a(e(y))) \rw e(y)$ is the whole third family as a schematic
rule, obtained by unification and rewriting, so it would need no proof by
induction.

\subsection{The braid monoid}

The three-strand positive braid monoid $\langle a, b \mid aba = bab\rangle$
has no finite convergent presentation over $\{a, b\}$~\cite{KapurNarendran85}.
With a Knuth--Bendix order where $a > b$, completion produces in the limit
the system $\{aba \rw bab\} \cup \{ab^nab \rw babba^{n-1} \mid n \ge 2\}$
\cite[Example~1]{HirokawaMiddeldorpSternagelWinkler17}. IsaScheme orients
the equation in the same direction and the critic recognizes exactly this
family, with exponent $n + 2$:
\begin{lstlisting}[language={}]
Divergence critic: a ((b ^^ Suc (Suc n)) (a (b x)))
                     = b (a (b (b ((a ^^ Suc n) x))))
Schematic rule not proved (built-in tactic)
\end{lstlisting}
The conjecture is a theorem, but the default tactic does not prove it
within 60\,s, so completion goes on diverging. The induction step needs
to rewrite $b(a(b(x)))$ to $a(b(a(x)))$ at the \emph{bottom} of the tower
$b^{n+3}$, while the definition of the function power exposes the copy of
$b$ at the \emph{top}. The standard lemma
$f\,(\fp{f}{n}\,x) = \fp{f}{n}\,(f\,x)$ (\isa{funpow\_swap1}) moves it
there, and either of the following declarations makes the proof go
through:
\begin{lstlisting}
declare funpow_swap1 [critic_simps]
(* or *)
divergence_critic_method
  ‹exponent_induct; (simp_all add: funpow_swap1)?; metis critic_rules›
\end{lstlisting}
With the first one, the critic proves the family in 0.2\,s, and completion
terminates after 2\,s with the system
\begin{gather*}
  a(b(a(x))) \rw b(a(b(x))), \\
  a\,(\fp{b}{\Suc\,(\Suc\;n)}\,(a\,(b\,x))) \rw
  b\,(a\,(b\,(b\,(\fp{a}{\Suc\;n}\,x)))).
\end{gather*}
The instances of the second rule, for $n \ge 0$, are exactly the rules
$ab^mab \rw babba^{m-1}$, $m \ge 2$, so this system has the same rewrite
relation as the infinite convergent system of the literature, and
completion has computed it as two rules over the original signature. (A finite
convergent system for the braid monoid needs an additional
symbol~\cite{Burckel01}.) With the second declaration the run that
finishes first is the one that orients the equation the other way,
$b(a(b(x))) \rw a(b(a(x)))$, and it ends after 3\,s with the mirror image
of this system.

\subsection{Associativity of Peano addition}

Dershowitz~\cite[Section~4.4]{Dershowitz89} uses the associativity of
Peano addition, with addition defined by recursion on its second argument,
to show that whether completion proves an inductive theorem depends on the
ordering. With one ordering completion succeeds, and with another it
generates the infinite sequence
\begin{gather*}
  (x+y)+z \approx x+(y+z),\qquad s(x+y)+z \approx x+(s(y)+z),\\
  s(s(x+y))+z \approx x+(s(s(y))+z),\quad \ldots
\end{gather*}
This example is closer to the intended use of IsaScheme, where completion
runs on the defining equations of recursive functions and on inductive
lemmas:
\begin{lstlisting}
datatype N = Z | S N
fun plus :: "N ⇒ N ⇒ N" where
  "plus x Z = x"
| "plus x (S y) = S (plus x y)"
declare [[generate_cps = true]]
complete "plus (plus x y) z = plus x (plus y z)"
  by (induction z, auto)
\end{lstlisting}
IsaScheme's completion orients associativity from left to right and runs
into Dershowitz's sequence. Along the way, its heuristic for equations of
the form $c(s) = c(t)$ proves the inductive lemma $\mathit{plus}\;Z\;z = z$,
and a second family appears. The critic generalizes both, in this order:
\begin{lstlisting}[language={}]
Schematic rule: plus ?x (plus ((S ^^ Suc ?n) Z) ?y)
                  ≡ plus ((S ^^ Suc ?n) ?x) ?y
Schematic rule: plus ((S ^^ Suc ?n) (plus ?x ?y)) ?z
                  ≡ plus ?x (plus ((S ^^ Suc ?n) ?y) ?z)
\end{lstlisting}
The second rule is exactly the family of Dershowitz's sequence. Completion
then terminates with the system
\begin{gather*}
  \mathit{plus}\;x\;Z \rw x, \qquad
  \mathit{plus}\;x\;(S\,y) \rw S\,(\mathit{plus}\;x\;y), \qquad
  \mathit{plus}\;Z\;z \rw z, \\
  \mathit{plus}\;(\mathit{plus}\;x\;y)\;z \rw \mathit{plus}\;x\;(\mathit{plus}\;y\;z), \\
  \mathit{plus}\;x\;(\mathit{plus}\;(\fp{S}{\Suc\;n}\,Z)\;y)
     \rw \mathit{plus}\;(\fp{S}{\Suc\;n}\,x)\;y, \\
  \mathit{plus}\;(\fp{S}{\Suc\;n}\,(\mathit{plus}\;x\;y))\;z
     \rw \mathit{plus}\;x\;(\mathit{plus}\;(\fp{S}{\Suc\;n}\,y)\;z)
\end{gather*}
in 73\,s. Both schematic rules are inductive theorems, proved by the
critic by induction on $n$ from the rules available at that point. Without
the critic, completion follows Dershowitz's sequence, together with the
second family, and has reached $n = 10$ when it is stopped after 300\,s.

\subsection{Families out of reach}

The last two examples diverge for reasons that the critic cannot address.
Completion of the theory of \emph{bands} (idempotent semigroups),
\[
  (x \cdot y) \cdot z \approx x \cdot (y \cdot z), \qquad x \cdot x \approx x,
\]
produces an infinite system with rules such as
$x \cdot (y \cdot (x \cdot y)) \rw x \cdot y$ and
$x \cdot (y \cdot (x \cdot (y \cdot z))) \rw x \cdot (y \cdot z)$
\cite[Section~2.4]{Dershowitz89}. Kirchner's \emph{signed binary
trees}~\cite[Example~2.1]{Kirchner89},
\begin{gather*}
  -(-x) \approx x,\qquad -f(x,y) \approx f(-y,-x),\\
  f(f(y,x),-x) \approx y,\qquad f(-x,f(x,y)) \approx y,
\end{gather*}
produce two families whose members grow as nested binary terms, such as
\begin{gather*}
  f(f(y,f(x_1,x_2)),f(-x_2,-x_1)) \rw y, \\
  f(f(y,f(x_1,f(x_2,x_3))),f(f(-x_3,-x_2),-x_1)) \rw y.
\end{gather*}
Kirchner schematizes them with meta-rules modulo the first two
equations~\cite{Kirchner89,BurghardtHeinz96}. In both theories the growing
part is a binary context, not a tower of a unary symbol, so the critic is
never triggered. IsaScheme reproduces the beginning of Kirchner's families,
and neither completion terminates within 300\,s. Representing these
families inside HOL needs a recursive function that builds the growing
context, which we discuss in Section~\ref{sec:conclusion}.

%% ===========================================================================
\section{Related work}
\label{sec:related}

\paragraph{Schematization of infinite rule sets.}
H.~Kirchner introduced \emph{meta-rules}, rules with meta-variables ranging
over terms, recognized by generalization modulo an equational theory, to
schematize the infinite families produced by divergent
completion~\cite{Kirchner87,Kirchner89}. Hermann characterized
\emph{crossed} rewrite systems, a syntactic source of
divergence~\cite{Hermann89}, and with Kirchner synthesized meta-rules from
them~\cite{KirchnerHermann90}. Chen, Hsiang and Kong proposed
\emph{recurrence terms}, terms with recursive definitions over a natural
number parameter, as finite representations of infinite sequences of
terms~\cite{ChenHsiangKong90}. Comon defined terms with integer exponents,
whose unification problem is decidable~\cite{Comon95}, and Hermann and
Galbav\'y generalized them to primal grammars~\cite{HermannGalbavy97}. Our
function powers are the special case of integer exponents applied to unary
contexts, which is also the case that Burghardt and Heinz identify as the
appropriate background theory for divergent examples like the one of
Example~\ref{ex:denzinger}~\cite{BurghardtHeinz96}. The difference is that
in HOL the exponent is part of the logic, so the schematic rule is a
theorem whose instances are obtained by instantiation and whose correctness
is proved by induction inside the prover, rather than argued about at the
meta level.

\paragraph{Recognizing divergence.}
Thomas and Jantke used inductive inference to replace an infinite sequence
of critical pairs by finitely many
generalizations~\cite{ThomasJantke89}, and Thomas and Watson replaced an
infinite rule set by a finite one over an enriched order-sorted
signature~\cite{ThomasWatson93}. Avenhaus achieved termination of
completion by introducing new function symbols and
equations~\cite{Avenhaus89}. Dershowitz and Pinchover used most specific
generalization to find divergence patterns among critical
pairs~\cite{DershowitzPinchover90}, and Burghardt and Heinz studied
anti-unification modulo equational theories as a tool for the same
purpose~\cite{BurghardtHeinz96}. Our detection is much simpler than
generalization modulo a theory, since it consists of exact matching of
skeletons plus an arithmetic check on tower lengths. It is cheap enough to run after every
choice of a rule, and it succeeds exactly on the families that a single
exponent can describe.

\paragraph{Divergence critics in inductive proof.}
Walsh's divergence critic~\cite{Walsh94,Walsh96} monitors an inductive
prover (Spike), recognizes accumulating term structure by difference
matching and speculates lemmas that ripple it away~\cite{Bundy05}. It is a
proof-planning critic in the sense of Ireland and Bundy~\cite{IrelandBundy96}.
Our critic shares its spirit, monitoring a procedure and proposing a lemma
that cuts the divergence. The lemma, however, is not a wave rule but the
generalization of the diverging sequence itself, which becomes expressible
thanks to the function power of HOL. IsaScheme's inductive
prover already includes a Walsh-style critic for its own proofs; the one
described here acts on completion.

\paragraph{Infinite and negative results.}
Some theories have no finite convergent presentation in their own
signature, the braid monoid $\langle a, b \mid aba = bab\rangle$ being the
classical example~\cite{KapurNarendran85}; Burckel gave a finite complete
system with an extra symbol~\cite{Burckel01}. Infinite complete systems can
still be useful. Hirokawa et al.\ formalized in Isabelle/HOL that fair
infinite runs produce complete presentations, with the braid monoid as
their running example~\cite{HirokawaMiddeldorpSternagelWinkler17}, and
the automatic groups programs of Epstein, Holt and Rees represent infinite
string rewriting systems by finite automata~\cite{EpsteinHoltRees91}.
Hermann and Pr\'{\i}vara studied the non-termination of the Knuth--Bendix
procedure~\cite{HermannPrivara86}; divergence depends on the reduction
ordering~\cite{Lescanne86,Dershowitz89}; and Sattler-Klein showed that the
infinite systems produced by completion can be very
complex~\cite{SattlerKlein91}. No critic based on a fixed class of
schematizations can therefore handle every divergent run.

\paragraph{Completion in proof assistants.}
Completion and Knuth--Bendix orders have been formalized in Isabelle/HOL as
part of the IsaFoR library~\cite{SternagelThiemann13,HirokawaMiddeldorpSternagel14,HirokawaMiddeldorpSternagelWinkler17}.
That work verifies completion as an algorithm; IsaScheme instead runs
completion \emph{inside} Isabelle on the user's theorems, and it is this
setting, where rules are theorems and induction is available, that makes
our generalization possible. Theory exploration systems such as
IsaCoSy~\cite{JohanssonDixonBundy11} and QuickSpec~\cite{ClaessenSmallboneHughes10}
face the same need for a non-redundant summary of the discovered equations.

%% ===========================================================================
\section{Limitations, future work and conclusion}
\label{sec:conclusion}

The critic is deliberately simple, and the experiments show where its
limits lie.

\paragraph{Shape of the families.}
Only towers of one unary symbol are recognized, all growing towers must
grow at the same rate, and the generalized rule has a single exponent.
Families such as $f(\pw{g}{2n}(x))$, families with two independent
exponents, and families that grow in non-unary contexts, like those of
bands and of Kirchner's signed binary trees, are out of reach. The first
two cases fit naturally in the same framework, since HOL can express
$g \pp (2 * n)$ and rules with several exponent
variables. The third one needs a recursive function that builds the
growing context. Inventing such functions is exactly what IsaScheme's
concept invention does with schemes~\cite{MontanoRivas12}, and connecting
the critic with it is the most interesting direction for future work.

\paragraph{Proving the generalization.}
The default tactic (induction, simplification and Metis with the current
rules) fails on the braid monoid although the family is a theorem, because
the induction step needs to move one copy of the iterated symbol from the
top to the bottom of the tower, and Metis does not find this step without
the lemma $f\,(\fp{f}{n}\,x) = \fp{f}{n}\,(f\,x)$. The configurable proof
method lets the user supply that lemma, or any other method, and then the
example is solved, but the user has to know what is missing. This could be
automated with a small default library of lemmas about the function power,
or with a method that calls IsaScheme's own inductive prover, which
analyses its failures and speculates lemmas. The second option is already
expressible
with \isa{divergence\_critic\_method}, but its cost has to be balanced
against the time limit of the critic.

\paragraph{Completeness of the restarted completion.}
The critical pairs of a schematic rule are approximated by those of its
first three instances, and the termination of the infinite family is
tested on the same instances. Both approximations could be replaced by
exact methods. The critical pairs of rules with exponents can be computed
by unification of terms with integer exponents~\cite{Comon95}, and the
termination of the schematic system could be proved with a reduction
ordering extended to function powers, or by proving the termination of the
instance $n$ by induction on $n$. With both, the final system would be convergent whenever
completion terminates.

\paragraph{Conclusion.}
In higher-order logic, the infinite families of rules produced by
divergent completion are often a single theorem about the function power.
We have shown how a theorem-producing completion procedure can detect such
families, state and prove the corresponding theorem by induction, and use
it to finish completion, without leaving the logic and without changing
the termination checks, the term orderings or the computation of critical
pairs. The
method solves the classical examples whose divergence is a growing tower of
a unary symbol, including Dershowitz's inductive example on Peano addition
and, with one lemma about the function power, the braid monoid, for which
it computes the known infinite convergent system as two rules. The examples
it does not solve point to natural extensions within the same framework.

%% ===========================================================================
\section*{Availability}
{\sloppy
IsaScheme/IsaLibs is available at
\url{https://github.com/omarmrivas/IsaScheme}. The theories used in
Section~\ref{sec:examples} are in the directory
\url{papers/divergent-completion-iterated-functions/experiments} of the
repository.\par}

\bibliographystyle{elsarticle-num}
\bibliography{references}

\end{document}
